\documentclass[11pt,a4paper]{report}

% ------------------------------------------------------------------
% Encodage & langue
% ------------------------------------------------------------------
\usepackage[utf8]{inputenc}
\usepackage[T1]{fontenc}
\usepackage[french]{babel}
\usepackage{lmodern}

% ------------------------------------------------------------------
% Mise en page
% ------------------------------------------------------------------
\usepackage[margin=2.3cm]{geometry}
\usepackage{graphicx}
\usepackage{xcolor}
\usepackage{listings}
\usepackage{booktabs}
\usepackage{array}
\usepackage{enumitem}
\usepackage{fancyhdr}
\usepackage{titlesec}
\usepackage[hidelinks]{hyperref}

% ------------------------------------------------------------------
% Couleurs
% ------------------------------------------------------------------
\definecolor{k8sblue}{RGB}{50,108,229}
\definecolor{codebg}{RGB}{245,245,245}
\definecolor{codekw}{RGB}{170,55,50}
\definecolor{codecomment}{RGB}{100,120,100}
\definecolor{codestr}{RGB}{60,120,60}

% ------------------------------------------------------------------
% En-tetes / pieds de page
% ------------------------------------------------------------------
\pagestyle{fancy}
\fancyhf{}
\lhead{\small Labs Kubernetes -- HPA \& Canary}
\rhead{\small ISI DITI4 2026}
\cfoot{\thepage}
\renewcommand{\headrulewidth}{0.4pt}

% ------------------------------------------------------------------
% Style des titres
% ------------------------------------------------------------------
\titleformat{\chapter}[display]{\huge\bfseries\color{k8sblue}}{\chaptertitlename\ \thechapter}{12pt}{\Huge}
\titlespacing*{\chapter}{0pt}{0pt}{20pt}
\titleformat{\section}{\Large\bfseries\color{k8sblue}}{\thesection}{0.6em}{}
\titleformat{\subsection}{\large\bfseries}{\thesubsection}{0.6em}{}

% ------------------------------------------------------------------
% Style du code
% ------------------------------------------------------------------
\lstdefinestyle{shell}{
  backgroundcolor=\color{codebg},
  basicstyle=\ttfamily\small,
  keywordstyle=\color{codekw}\bfseries,
  commentstyle=\color{codecomment}\itshape,
  stringstyle=\color{codestr},
  breaklines=true,
  frame=single,
  rulecolor=\color{gray!40},
  showstringspaces=false,
  columns=fullflexible,
  keepspaces=true,
  xleftmargin=6pt,
  framexleftmargin=6pt,
}
\lstdefinelanguage{yaml}{
  keywords={apiVersion,kind,metadata,spec,name,namespace,labels,replicas,selector,matchLabels,template,containers,image,ports,resources,requests,limits,cpu,memory,strategy,type,rollingUpdate,maxSurge,maxUnavailable,minReplicas,maxReplicas,metrics,resource,target,averageUtilization,scaleTargetRef,hard,pods,ports,targetPort},
  keywordstyle=\color{k8sblue}\bfseries,
  comment=[l]{\#},
  commentstyle=\color{codecomment}\itshape,
  morestring=[b]',
  stringstyle=\color{codestr},
  basicstyle=\ttfamily\small,
  backgroundcolor=\color{codebg},
  breaklines=true,
  frame=single,
  rulecolor=\color{gray!40},
  showstringspaces=false,
  xleftmargin=6pt,
  framexleftmargin=6pt,
}
\lstset{style=shell}

% Commande pour inserer une capture d'ecran (image presente)
\newcommand{\capture}[2]{%
  \begin{figure}[h!]
    \centering
    \includegraphics[width=0.92\textwidth]{captures/#2}
    \caption{#1}
    \label{fig:#2}
  \end{figure}%
}
% Emplacement pour une capture pas encore prise (n'insere pas d'image)
\newcommand{\captureTODO}[2]{%
  \begin{figure}[h!]
    \centering
    \fbox{\parbox[c][4cm][c]{0.85\textwidth}{\centering\itshape
      \textcolor{gray}{[ Capture a inserer : #1 ]}\\[4pt]
      \small\texttt{captures/#2.png}}}
    \caption{#1}
    \label{fig:#2}
  \end{figure}%
}

% ==================================================================
\begin{document}

% ---------- Page de titre ----------
\begin{titlepage}
  \centering
  \vspace*{2cm}
  {\Huge\bfseries\color{k8sblue} Laboratoires Kubernetes\par}
  \vspace{0.6cm}
  {\Large Autoscaling (HPA) \& Canary Release\par}
  \vspace{1cm}
  {\large Deploiement, mise a l'echelle automatique et livraison progressive\\[4pt]
   d'une application web sur un cluster local\par}
  \vspace{0.4cm}
  {\large Docker Desktop -- Prometheus / Grafana -- K6\par}
  \vfill
  \begin{tabular}{rl}
    \textbf{Formation} & : ISI -- DITI4 2026 \\
    \textbf{Module}    & : DevOps / Cloud \\
    \textbf{Etudiant}  & : \underline{\hspace{5cm}} \\
    \textbf{Date}      & : \today \\
  \end{tabular}
  \vfill
\end{titlepage}

\tableofcontents
\newpage

% ##################################################################
% CHAPITRE 1
% ##################################################################
\chapter{Mise a l'echelle automatique (HPA)}

% ==================================================================
\section{Objectif du laboratoire}

L'objectif de ce laboratoire est de mettre en place et de \textbf{demontrer la mise a
l'echelle automatique horizontale} (\emph{Horizontal Pod Autoscaling}, HPA) d'une
application web sur un cluster Kubernetes.

\medskip
Concretement, il s'agit de prouver que Kubernetes est capable :
\begin{itemize}[nosep]
  \item d'\textbf{ajouter automatiquement des copies} (pods) de l'application lorsque la
        charge (CPU) augmente ;
  \item de \textbf{reduire automatiquement} le nombre de pods lorsque la charge diminue.
\end{itemize}

\begin{quote}
\textbf{Analogie.} Une boulangerie garde \emph{un} vendeur en temps normal. A l'heure de
pointe, le patron appelle des vendeurs supplementaires ; quand la file se vide, il les
renvoie. Le HPA joue le role de ce patron automatique, en se basant sur la consommation
CPU des pods.
\end{quote}

% ==================================================================
\section{Architecture et composants}

\begin{table}[h!]
\centering
\renewcommand{\arraystretch}{1.3}
\begin{tabular}{>{\bfseries}p{3.2cm} p{10.5cm}}
\toprule
Composant & Role dans le laboratoire \\
\midrule
Deployment \texttt{webapp} & L'application a mettre a l'echelle (1 pod au depart). \\
Service \texttt{webapp}    & Expose l'application a l'interieur du cluster (port 80). \\
metrics-server             & Mesure la consommation CPU/memoire des pods (le << thermometre >>). Indispensable au HPA. \\
HPA \texttt{myapp-hpa}     & Regle d'autoscaling : si CPU moyen $>$ 50\%, ajouter des pods (max 10). \\
ResourceQuota              & Limite les ressources totales du namespace. \\
Prometheus / Grafana       & Collecte et visualisation des metriques (courbes). \\
K6                         & Generateur de charge : simule 1000 utilisateurs pour declencher l'autoscaling. \\
\bottomrule
\end{tabular}
\caption{Composants deployes dans le namespace \texttt{lab-hpa}.}
\end{table}

% ==================================================================
\section{Prerequis / Environnement}

\begin{lstlisting}[language=bash]
kubectl version --client   # v1.32.2
kubectl config current-context   # docker-desktop
helm version               # v4.1.4
docker version             # 28.3.0
kubectl get nodes          # desktop-control-plane   Ready
\end{lstlisting}

Le cluster utilise est celui integre a \textbf{Docker Desktop} (un noeud
\texttt{desktop-control-plane}). Toutes les commandes sont executees depuis
\textbf{PowerShell}.

% ==================================================================
\section{Etape 1 -- Creation du namespace}

Un \emph{namespace} isole les ressources du laboratoire du reste du cluster.

\begin{lstlisting}[language=bash]
kubectl create ns lab-hpa
\end{lstlisting}

% ==================================================================
\section{Etape 2 -- Deploiement de l'application}

Le fichier \texttt{deploy-lab-hpa.yaml} decrit le deploiement de l'application
(image \texttt{aichalo/intranet-newdeal:latest}). Les sections \texttt{requests} et
\texttt{limits} sont \textbf{obligatoires} : sans \texttt{requests.cpu}, le HPA ne peut
pas calculer le pourcentage d'utilisation CPU.

\begin{lstlisting}[language=yaml]
apiVersion: apps/v1
kind: Deployment
metadata:
  name: webapp
  namespace: lab-hpa
  labels:
    app: webapp
spec:
  replicas: 1
  selector:
    matchLabels:
      app: webapp
  strategy:
    type: RollingUpdate
    rollingUpdate:
      maxSurge: 25%
      maxUnavailable: 25%
  template:
    metadata:
      labels:
        app: webapp
    spec:
      containers:
        - name: webapp
          image: aichalo/intranet-newdeal:latest
          imagePullPolicy: IfNotPresent
          ports:
            - containerPort: 80
          resources:
            requests:
              cpu: 20m
              memory: 56Mi
            limits:
              cpu: 20m
              memory: 56Mi
\end{lstlisting}

Le fichier \texttt{service.yaml} expose l'application (necessaire pour les tests K6) :

\begin{lstlisting}[language=yaml]
apiVersion: v1
kind: Service
metadata:
  name: webapp
  namespace: lab-hpa
  labels:
    app: webapp
spec:
  type: ClusterIP
  selector:
    app: webapp
  ports:
    - name: http
      port: 80
      targetPort: 80
\end{lstlisting}

\begin{lstlisting}[language=bash]
kubectl apply -f deploy-lab-hpa.yaml
kubectl apply -f service.yaml
kubectl get pods -n lab-hpa
\end{lstlisting}

% ==================================================================
\section{Etape 3 -- Stack d'observabilite (Prometheus / Grafana / Loki)}

\begin{lstlisting}[language=bash]
kubectl create namespace observability

helm repo add prometheus-community https://prometheus-community.github.io/helm-charts
helm repo add grafana https://grafana.github.io/helm-charts
helm repo update

helm install monitoring prometheus-community/kube-prometheus-stack -n observability
helm install loki grafana/loki-stack -n observability

kubectl get pods -n observability
\end{lstlisting}

Acces a Grafana (dans un terminal dedie) :

\begin{lstlisting}[language=bash]
kubectl port-forward svc/monitoring-grafana 8085:80 -n observability
# Mot de passe admin :
$p = kubectl get secret monitoring-grafana -n observability -o jsonpath="{.data.admin-password}"
[System.Text.Encoding]::UTF8.GetString([System.Convert]::FromBase64String($p))
\end{lstlisting}

Interface accessible sur \url{http://localhost:8085} (utilisateur \texttt{admin}).

% ==================================================================
\section{Etape 4 -- Installation du metrics-server}

\textbf{Correction par rapport au sujet :} l'URL du depot Helm du sujet
(\texttt{https://github.io}) est erronee. La bonne URL est
\texttt{https://kubernetes-sigs.github.io/metrics-server/}. Le drapeau
\texttt{-{}-kubelet-insecure-tls} est necessaire sur Docker Desktop.

\begin{lstlisting}[language=bash]
helm repo add metrics-server https://kubernetes-sigs.github.io/metrics-server/
helm repo update
helm install metrics-server metrics-server/metrics-server -n kube-system \
  --set "args={--kubelet-insecure-tls}"

# Verification
kubectl get pods -n kube-system -l app.kubernetes.io/name=metrics-server
kubectl top nodes
kubectl top pods -n lab-hpa
\end{lstlisting}

% ==================================================================
\section{Etape 5 -- ResourceQuota}

\begin{lstlisting}[language=yaml]
apiVersion: v1
kind: ResourceQuota
metadata:
  name: app-quota
  namespace: lab-hpa
spec:
  hard:
    pods: "50"
    requests.cpu: "500m"
    requests.memory: "1Gi"
    limits.cpu: "2000m"
    limits.memory: "2Gi"
\end{lstlisting}

\begin{lstlisting}[language=bash]
kubectl apply -f quota-app.yaml
kubectl get resourcequota -n lab-hpa
\end{lstlisting}

% ==================================================================
\section{Etape 6 -- Deploiement du HPA}

\textbf{Correction par rapport au sujet :} ajout de \texttt{namespace: lab-hpa} (absent du
sujet), sinon le HPA cible le namespace \texttt{default}.

\begin{lstlisting}[language=yaml]
apiVersion: autoscaling/v2
kind: HorizontalPodAutoscaler
metadata:
  name: myapp-hpa
  namespace: lab-hpa
spec:
  scaleTargetRef:
    apiVersion: apps/v1
    kind: Deployment
    name: webapp
  minReplicas: 1
  maxReplicas: 10
  metrics:
    - type: Resource
      resource:
        name: cpu
        target:
          type: Utilization
          averageUtilization: 50
    - type: Resource
      resource:
        name: memory
        target:
          type: Utilization
          averageUtilization: 50
\end{lstlisting}

\begin{lstlisting}[language=bash]
kubectl apply -f hpa.yaml
kubectl get hpa -n lab-hpa
\end{lstlisting}

Au repos, le HPA affiche les cibles avec des pourcentages et 1 replica :

\begin{lstlisting}[language=bash]
NAME        REFERENCE           TARGETS                        MINPODS  MAXPODS  REPLICAS
myapp-hpa   Deployment/webapp   cpu: 0%/50%, memory: 17%/50%   1        10       1
\end{lstlisting}

\capture{HPA au repos (\texttt{cpu: 0\%/50\%}, REPLICAS 1) puis, apres lancement de la charge, \texttt{cpu: 99\%/50\%} et REPLICAS 10}{04-hpa-repos}

% ==================================================================
\section{Etape 7 -- Test de charge avec K6}

On expose d'abord l'application sur le port 8087 (terminal dedie) :

\begin{lstlisting}[language=bash]
kubectl port-forward svc/webapp 8087:80 -n lab-hpa
\end{lstlisting}

\subsection{Scenario 1 -- charge constante}

\begin{lstlisting}[language=bash]
# script.js : 1000 utilisateurs virtuels pendant 6 minutes
docker run --rm -i grafana/k6 run - < script.js
\end{lstlisting}

\subsection{Scenario 2 -- montee/descente progressive}

\begin{lstlisting}[language=bash]
# ramp-up-down.js : 1000 -> 500 -> 250 -> 0 utilisateurs
docker run --rm -i grafana/k6 run - < ramp-up-down.js
\end{lstlisting}

\subsection{Scenario retenu -- generateur de charge interne au cluster}

\textbf{Constat.} Lancer K6 depuis l'exterieur via \texttt{kubectl port-forward} n'a pas
suffi a faire monter le CPU : le \texttt{port-forward} constitue un goulot d'etranglement
(flux limite) et l'application (nginx statique) consomme tres peu de CPU. Les requetes
n'atteignaient donc pas les pods en quantite suffisante.

\medskip
\textbf{Solution.} Generer la charge \emph{depuis l'interieur} du cluster, directement sur
le Service \texttt{webapp} (ClusterIP), avec un deploiement dedie de 4 pods executant
chacun 10 boucles \texttt{wget} en parallele. Comme le \texttt{ResourceQuota} impose que
tout pod declare ses \texttt{requests}/\texttt{limits}, ceux-ci sont inclus.

\begin{lstlisting}[language=yaml]
apiVersion: apps/v1
kind: Deployment
metadata:
  name: load-generator
  namespace: lab-hpa
spec:
  replicas: 4
  selector:
    matchLabels:
      app: load-generator
  template:
    metadata:
      labels:
        app: load-generator
    spec:
      containers:
        - name: load
          image: busybox:1.28
          command: ["/bin/sh", "-c"]
          args:
            - >
              i=0; while [ $i -lt 10 ]; do
                ( while true; do wget -q -O- http://webapp.lab-hpa.svc.cluster.local >/dev/null 2>&1; done ) &
                i=$((i+1)); done; wait
          resources:
            requests: { cpu: 50m, memory: 16Mi }
            limits:   { cpu: 200m, memory: 32Mi }
\end{lstlisting}

\begin{lstlisting}[language=bash]
kubectl apply -f load-generator.yaml     # demarre la charge
kubectl delete -f load-generator.yaml    # arrete la charge (declenche le scale-down)
\end{lstlisting}

Pendant les tests, on observe la mise a l'echelle :

\begin{lstlisting}[language=bash]
kubectl get pods -n lab-hpa -w
kubectl get hpa -n lab-hpa
\end{lstlisting}

\capture{Montee en charge (terminal) : \texttt{cpu: 99\%/50\%}, REPLICAS 10 et 10 pods \texttt{webapp} crees automatiquement}{05-scaling-up}
\capture{Montee en charge (Grafana) : CPU Utilisation 246\% (from requests), courbe CPU au plafond du quota}{06-grafana-scaling}
\capture{Retour au calme (terminal) : le HPA supprime progressivement les pods jusqu'a 1 seul \texttt{webapp}}{07-scaling-down}
\capture{Retour au calme (Grafana) : CPU Utilisation retombee a 0\% apres l'arret de la charge}{08-grafana-down}

% ==================================================================
\section{Observations et resultats}

\begin{table}[h!]
\centering
\renewcommand{\arraystretch}{1.3}
\begin{tabular}{p{4.5cm} p{4cm} p{4cm}}
\toprule
\textbf{Phase} & \textbf{CPU (cible 50\%)} & \textbf{Nombre de pods webapp} \\
\midrule
Repos (avant charge)        & 0\%            & 1 \\
Montee en charge            & $>$ 50\%       & 1 $\rightarrow$ 10 \\
Pic de charge               & 99\%           & 10 (maximum atteint) \\
Apres arret de la charge    & retombe a 0\%  & 10 $\rightarrow$ 1 \\
\bottomrule
\end{tabular}
\caption{Evolution mesuree du nombre de pods en fonction de la charge CPU.}
\end{table}

\noindent Sortie du HPA au pic de charge (\texttt{kubectl get hpa -n lab-hpa}) :

\begin{lstlisting}[language=bash]
NAME        REFERENCE           TARGETS                         MINPODS  MAXPODS  REPLICAS
myapp-hpa   Deployment/webapp   cpu: 99%/50%, memory: 18%/50%   1        10       10
\end{lstlisting}

\noindent Lecture Grafana au pic (namespace \texttt{lab-hpa}) :
\begin{itemize}[nosep]
  \item CPU Utilisation (from requests) : \textbf{246\%}
  \item CPU Utilisation (from limits)   : 98,2\%
  \item Memory Utilisation (from requests) : 16,9\%
\end{itemize}

\textbf{Conclusion.} Le laboratoire demontre que le HPA, couple au metrics-server,
ajuste automatiquement le nombre de pods de l'application en fonction de la charge CPU.
Sous l'effet du generateur de charge, le CPU moyen des pods \texttt{webapp} est passe de
0\% a 99\% (cible : 50\%), ce qui a declenche la creation automatique de pods
supplementaires jusqu'a la limite fixee (\texttt{maxReplicas: 10}). Apres l'arret de la
charge, le CPU est retombe et le HPA a ramene le deploiement a 1 pod (apres la fenetre de
stabilisation de 5 minutes), confirmant le fonctionnement complet de la mise a l'echelle
automatique (montee \emph{et} descente).

% ==================================================================
\section{Nettoyage}

\begin{lstlisting}[language=bash]
helm uninstall monitoring -n observability
helm uninstall loki -n observability
kubectl delete namespace observability
kubectl delete namespace lab-hpa
helm uninstall metrics-server -n kube-system
\end{lstlisting}

% ==================================================================
\section{Difficultes rencontrees et solutions}

\begin{table}[h!]
\centering
\renewcommand{\arraystretch}{1.5}
\begin{tabular}{p{5.2cm} p{8.3cm}}
\toprule
\textbf{Probleme} & \textbf{Cause et solution} \\
\midrule
Le CPU ne monte pas malgre le test K6, le HPA reste a 1 replica. &
K6 lance depuis l'exterieur via \texttt{kubectl port-forward} est bride (un seul flux) et
l'application nginx est tres legere. Solution : generer la charge \emph{depuis l'interieur}
du cluster (deploiement \texttt{load-generator} tapant sur le Service ClusterIP). \\
\addlinespace
\texttt{base64 : n'est pas reconnu} lors du decodage du mot de passe Grafana. &
La commande \texttt{base64} n'existe pas sous Windows. Solution : decoder en PowerShell avec
\texttt{[System.Convert]::FromBase64String(\$p)}. \\
\addlinespace
\texttt{kubectl get hpa,pods -w} : \texttt{you may only specify a single resource type}. &
Le mode \emph{watch} (\texttt{-w}) n'accepte qu'un seul type de ressource. Solution :
surveiller les pods et le HPA dans deux commandes distinctes. \\
\addlinespace
\texttt{kubectl top nodes} : \texttt{Metrics API not available}. &
Le metrics-server venait de demarrer et n'etait pas encore \texttt{Ready}. Solution :
attendre $\sim$1 min que son pod soit \texttt{1/1 Running}, puis relancer. \\
\addlinespace
Installation du metrics-server impossible (depot Helm introuvable). &
L'URL du sujet (\texttt{https://github.io}) est erronee. Solution : utiliser
\texttt{https://kubernetes-sigs.github.io/metrics-server/} et ajouter l'argument
\texttt{-{}-kubelet-insecure-tls} (obligatoire sur Docker Desktop). \\
\bottomrule
\end{tabular}
\caption{Journal des difficultes rencontrees et de leurs solutions (chapitre HPA).}
\end{table}

\section*{Annexe -- Corrections apportees au sujet}
\addcontentsline{toc}{section}{Annexe -- Corrections apportees au sujet}

\begin{enumerate}[nosep]
  \item Manifeste \texttt{Deployment} reformate proprement (le YAML \texttt{items:} du sujet
        etait mal indente) et utilisation de l'image du projet \texttt{aichalo/intranet-newdeal}.
  \item Ajout d'un \texttt{Service} \texttt{webapp} (absent du sujet mais requis pour le
        port-forward vers le port 8087 utilise par K6).
  \item metrics-server : URL du depot corrigee en
        \texttt{https://kubernetes-sigs.github.io/metrics-server/} et ajout de
        \texttt{-{}-kubelet-insecure-tls}.
  \item Loki : \texttt{port-forward} sur \texttt{-n observability} (au lieu de
        \texttt{-n monitoring}).
  \item HPA : ajout de \texttt{namespace: lab-hpa}.
\end{enumerate}

% ##################################################################
% CHAPITRE 2
% ##################################################################
\chapter{Canary Release}

\section{Objectif}

Le \emph{Canary Release} (livraison << canari >>) consiste a deployer une nouvelle version
d'une application (v2) a cote de la version en production (v1), en n'envoyant d'abord
qu'une \textbf{petite part du trafic} vers v2. Si la v2 se comporte correctement, on
augmente progressivement sa part (20\%, 50\%, 100\%) ; en cas de probleme, on revient a
v1 sans impact notable pour les utilisateurs.

\begin{quote}
\textbf{Analogie.} On envoie un << canari dans la mine >> : une petite fraction des
utilisateurs teste la nouvelle version. Si tout va bien, on generalise ; sinon, on fait
marche arriere.
\end{quote}

\section{Methode retenue : Kubernetes natif}

Aucun outil supplementaire (ni Istio, ni Argo Rollouts). Le principe repose sur un
\textbf{unique Service} dont le selecteur est volontairement << large >> (\texttt{app:
intranet}) et qui cible donc \textbf{deux Deployments} partageant ce label mais
differencies par \texttt{version: v1} / \texttt{version: v2}.

\begin{quote}
Le Service repartit le trafic (round-robin) sur tous les pods correspondants. Par
consequent : \textbf{la part de trafic vers v2 = proportion de pods v2 sur le total}.
\end{quote}

\section{Preparation de la version v2}

La v2 est identique a la v1, a l'exception d'un bandeau visible << VERSION 2 >> ajoute dans
\texttt{index.html} (marqueur permettant de distinguer v1 et v2 a l'oeil et au \texttt{grep}).
L'image v2 repart de la v1 et ne remplace que ce fichier :

\begin{lstlisting}[language=bash]
# lab-canary/v2/Dockerfile
FROM aichalo/intranet-newdeal:latest
COPY index.html /usr/share/nginx/html/index.html
\end{lstlisting}

\begin{lstlisting}[language=bash]
cd lab-canary/v2
docker build -t aichalo/intranet-newdeal:2.0 .
docker push aichalo/intranet-newdeal:2.0
\end{lstlisting}

\capture{Site en version v1 (bandeau normal)}{c2-01-browser-v1}
\capture{Site en version v2 (bandeau rouge << VERSION 2 >>, titre suivi de (v2))}{c2-02-browser-v2}

\section{Manifests}

Deux Deployments et un Service. Extrait des labels (le reste : image, ports, ressources,
readinessProbe) :

\begin{lstlisting}[language=yaml]
# deploy-v1.yaml (extrait)      # deploy-v2.yaml (extrait)
metadata:                       metadata:
  name: intranet-v1              name: intranet-v2
spec:                           spec:
  replicas: 4                    replicas: 0     # aucun trafic v2 au depart
  template:                      template:
    metadata:                      metadata:
      labels:                        labels:
        app: intranet                app: intranet   # label COMMUN
        version: v1                  version: v2
    ... image: :latest            ... image: :2.0
\end{lstlisting}

\begin{lstlisting}[language=yaml]
# service.yaml -- LE point cle : selecteur SANS "version"
apiVersion: v1
kind: Service
metadata:
  name: intranet
  namespace: lab-canary
spec:
  selector:
    app: intranet        # cible v1 ET v2
  ports:
    - port: 80
      targetPort: 80
\end{lstlisting}

\begin{lstlisting}[language=bash]
kubectl create ns lab-canary
kubectl apply -f deploy-v1.yaml
kubectl apply -f deploy-v2.yaml
kubectl apply -f service.yaml
kubectl get pods -n lab-canary
\end{lstlisting}

\capture{Etat initial : 4 pods \texttt{intranet-v1} Running (100\% v1, 0 pod v2)}{c2-06-pods-init}

\section{Methode de mesure de la repartition}

\textbf{Point technique important.} Le Kubernetes integre a Docker Desktop (noeud
\texttt{desktop-control-plane}, base sur \emph{kind}) \textbf{n'expose pas les Services
NodePort sur \texttt{localhost}} ; de plus, \texttt{kubectl port-forward} se connecte a
\emph{un seul} pod et ne repartit donc pas le trafic. Ces deux methodes ne permettent pas
d'observer le partage v1/v2.

\medskip
La repartition n'est reelle que pour le trafic \textbf{interne au cluster}, ou kube-proxy
load-balance le Service (ClusterIP) entre tous les endpoints. On mesure donc depuis
l'interieur du cluster : un pod \texttt{busybox} interroge $N$ fois le Service via son nom
DNS interne et compte les reponses de la v2 (reperees par l'identifiant \texttt{canary-banner},
present uniquement dans son \texttt{index.html}). C'est le role du script
\texttt{count-versions.ps1} :

\begin{lstlisting}[language=bash]
kubectl run canary-count --rm -i --restart=Never -n lab-canary \
  --image=busybox:1.28 --command -- sh -c \
  'n=50; v2=0; i=0; while [ $i -lt $n ]; do \
     wget -q -O - http://intranet.lab-canary.svc.cluster.local \
       | grep -q canary-banner && v2=$((v2+1)); \
     i=$((i+1)); done; echo RESULT v1=$((n-v2)) v2=$v2'
\end{lstlisting}

\section{Deroule de la progression}

Pour chaque palier, on ajuste le nombre de replicas, on attend que les pods soient
\texttt{Running}, puis on lance \texttt{count-versions.ps1}.

\begin{table}[h!]
\centering
\renewcommand{\arraystretch}{1.4}
\begin{tabular}{p{2.5cm} p{7.5cm} p{2.5cm}}
\toprule
\textbf{Palier} & \textbf{Commande} & \textbf{\% trafic v2} \\
\midrule
Initial   & v1=4, v2=0 (etat de depart) & 0\% \\
Canary    & \texttt{kubectl scale deploy/intranet-v2 --replicas=1} & $\approx$20\% \\
Equilibre & \texttt{scale intranet-v1 --replicas=2} et \texttt{intranet-v2 --replicas=2} & 50\% \\
Bascule   & \texttt{scale intranet-v1 --replicas=0} et \texttt{intranet-v2 --replicas=4} & 100\% \\
\bottomrule
\end{tabular}
\caption{Progression du canary pilotee par le nombre de pods.}
\end{table}

\begin{lstlisting}[language=bash]
# Exemple : passer au canary 20%
kubectl scale deploy/intranet-v2 -n lab-canary --replicas=1
.\count-versions.ps1
\end{lstlisting}

\capture{Palier canary 20\% : 4 pods v1 + 1 pod v2, comptage \texttt{v1=37 v2=13}}{c2-03-count-20}
\capture{Palier 50/50 : 2 pods v1 + 2 pods v2, comptage \texttt{v1=27 v2=23}}{c2-04-count-50}
\capture{Palier 100\% v2 : 4 pods v2, comptage \texttt{v1=0 v2=50}}{c2-05-count-100}

\section{Rollback (securite du canary)}

L'un des interets majeurs du canary est de pouvoir \textbf{annuler instantanement} en cas
de probleme sur v2, en ramenant tout le trafic vers v1 :

\begin{lstlisting}[language=bash]
kubectl scale deploy/intranet-v1 -n lab-canary --replicas=4
kubectl scale deploy/intranet-v2 -n lab-canary --replicas=0
\end{lstlisting}

\section{Difficultes rencontrees et solutions}

Le tableau ci-dessous recense les problemes reels rencontres pendant la realisation du
laboratoire et la maniere dont ils ont ete resolus.

\begin{table}[h!]
\centering
\renewcommand{\arraystretch}{1.5}
\begin{tabular}{p{5.2cm} p{8.3cm}}
\toprule
\textbf{Probleme} & \textbf{Cause et solution} \\
\midrule
Le site est inaccessible sur \texttt{localhost:30090} (\texttt{curl} renvoie HTTP 000). &
Docker Desktop (Kubernetes type \emph{kind}) n'expose pas les NodePort sur l'hote. Solution :
mesurer la repartition depuis l'interieur du cluster (Service ClusterIP + kube-proxy). \\
\addlinespace
Le comptage donne toujours 100\% v1, quel que soit le nombre de pods v2. &
\texttt{kubectl port-forward svc/...} se connecte a un \emph{seul} pod et ne repartit pas le
trafic. Solution : interroger le Service via son DNS interne depuis un pod \texttt{busybox}. \\
\addlinespace
Le script de comptage plante : \texttt{syntax error: unexpected end of file (expecting "then")}. &
Le marqueur recherche \texttt{"VERSION 2"} contient un espace ; les guillemets etaient mal
echappes lors du passage PowerShell $\rightarrow$ kubectl $\rightarrow$ shell. Solution :
chercher un marqueur sans espace ni guillemet (\texttt{canary-banner}) et passer le script
shell en quotes simples cote PowerShell. \\
\addlinespace
Requetes \texttt{wget: can't connect} pendant un comptage. &
Le comptage a ete lance alors que les pods etaient encore en \texttt{ContainerCreating} /
\texttt{Terminating}. Solution : attendre que tous les pods soient \texttt{1/1 Running}
avant de mesurer. \\
\addlinespace
\texttt{.\textbackslash count-versions.ps1 : n'est pas reconnu}. &
La commande a ete lancee hors du dossier \texttt{lab-canary}. Solution : se placer dans le
dossier (\texttt{cd}) avant d'executer le script. \\
\bottomrule
\end{tabular}
\caption{Journal des difficultes rencontrees et de leurs solutions.}
\end{table}

\section{Resultats et conclusion}

\begin{table}[h!]
\centering
\renewcommand{\arraystretch}{1.3}
\begin{tabular}{p{2.6cm} p{3cm} p{3.2cm} p{3cm}}
\toprule
\textbf{Palier} & \textbf{Ratio pods (v1/v2)} & \textbf{Comptage (sur 50)} & \textbf{Trafic v2 mesure} \\
\midrule
Initial   & 4 / 0 & v1=50, v2=0  & 0\% \\
Canary    & 4 / 1 & v1=37, v2=13 & 26\% \\
Equilibre & 2 / 2 & v1=27, v2=23 & 46\% \\
Bascule   & 0 / 4 & v1=0,  v2=50 & 100\% \\
\bottomrule
\end{tabular}
\caption{Repartition mesuree du trafic a chaque palier (comptage in-cluster sur 50 requetes).}
\end{table}

\noindent Les proportions mesurees (0\%, 26\%, 46\%, 100\%) correspondent bien aux ratios
de pods configures (0/5, 1/5, 2/4, 4/4). Les legers ecarts (26\% au lieu de 20\%, 46\% au
lieu de 50\%) proviennent de la variabilite normale du round-robin sur un faible nombre de
requetes.

\textbf{Conclusion.} Le laboratoire demontre une strategie de \emph{Canary Release} en
Kubernetes natif : grace a un Service ciblant un label commun a deux Deployments, la part
de trafic dirigee vers la nouvelle version est controlee par le nombre de pods. La montee
en charge progressive (0\% $\rightarrow$ 20\% $\rightarrow$ 50\% $\rightarrow$ 100\%) et la
possibilite de \emph{rollback} immediat illustrent l'interet du canary pour livrer une
nouvelle version en limitant le risque.

\end{document}
